\documentclass[10pt,a4paper]{amsart}
\usepackage{amsmath,amssymb,amsthm,graphicx}
% For commutative diagrams we use Xy-pic
\usepackage[all]{xy}

% We use 2cell for 2-commutative diagrams.
\xyoption{2cell}
% Theorem environments.
%
\theoremstyle{plain}
\newtheorem{theorem}[subsection]{Theorem}
\newtheorem{proposition}[subsection]{Proposition}
\newtheorem{lemma}[subsection]{Lemma}

\theoremstyle{definition}
\newtheorem{definition}[subsection]{Definition}
\newtheorem{example}[subsection]{Example}
\newtheorem{exercise}[subsection]{Exercise}
\newtheorem{situation}[subsection]{Situation}

\theoremstyle{remark}
\newtheorem{remark}[subsection]{Remark}
\newtheorem{remarks}[subsection]{Remarks}

\numberwithin{equation}{subsection}

% Macros
%
\def\lim{\mathop{\mathrm{lim}}\nolimits}
\def\colim{\mathop{\mathrm{colim}}\nolimits}
\def\Spec{\mathop{\mathrm{Spec}}}
\def\Hom{\mathop{\mathrm{Hom}}\nolimits}
\def\Ext{\mathop{\mathrm{Ext}}\nolimits}
\def\SheafHom{\mathop{\mathcal{H}\!\mathit{om}}\nolimits}
\def\SheafExt{\mathop{\mathcal{E}\!\mathit{xt}}\nolimits}
\def\Sch{\mathit{Sch}}
\def\Mor{\mathop{\mathrm{Mor}}\nolimits}
\def\Ob{\mathop{\mathrm{Ob}}\nolimits}
\def\Sh{\mathop{\mathit{Sh}}\nolimits}
\def\NL{\mathop{N\!L}\nolimits}
\def\CH{\mathop{\mathrm{CH}}\nolimits}
\def\proetale{{pro\text{-}\acute{e}tale}}
\def\etale{{\acute{e}tale}}
\def\QCoh{\mathit{QCoh}}
\def\Ker{\mathop{\mathrm{Ker}}}
\def\Im{\mathop{\mathrm{Im}}}
\def\Coker{\mathop{\mathrm{Coker}}}
\def\Coim{\mathop{\mathrm{Coim}}}

% Boxtimes
%
\DeclareMathSymbol{\boxtimes}{\mathbin}{AMSa}{"02}

%
% Macros for moduli stacks/spaces
%
\def\QCohstack{\mathcal{QC}\!\mathit{oh}}
\def\Cohstack{\mathcal{C}\!\mathit{oh}}
\def\Spacesstack{\mathcal{S}\!\mathit{paces}}
\def\Quotfunctor{\mathrm{Quot}}
\def\Hilbfunctor{\mathrm{Hilb}}
\def\Curvesstack{\mathcal{C}\!\mathit{urves}}
\def\Polarizedstack{\mathcal{P}\!\mathit{olarized}}
\def\Complexesstack{\mathcal{C}\!\mathit{omplexes}}
% \Pic is the operator that assigns to X its picard group, usage \Pic(X)
% \Picardstack_{X/B} denotes the Picard stack of X over B
% \Picardfunctor_{X/B} denotes the Picard functor of X over B
\def\Pic{\mathop{\mathrm{Pic}}\nolimits}
\def\Picardstack{\mathcal{P}\!\mathit{ic}}
\def\Picardfunctor{\mathrm{Pic}}
\def\Deformationcategory{\mathcal{D}\!\mathit{ef}}
\newenvironment{reference}{}{}
\newenvironment{slogan}{}{}
\newenvironment{history}{}{}
\usepackage{cleveref}
\crefname{originalsection}{original section}{original sections}
\Crefname{originalsection}{Original section}{Original sections}
\makeatletter
\newlabel{discriminant-section-tate-map}{{0FKB}{source}{}{sourcepointer-1}{}}
\newlabel{discriminant-section-tate-map@cref}{{[originalsection][0][]0FKB}{[1][0][]source}{}{}{}}
\makeatother
\crefname{originalsection}{original section}{original sections}
\Crefname{originalsection}{Original section}{Original sections}
\makeatletter
\newlabel{section-trace}{{0FK6}{source}{}{sourcepointer-2}{}}
\newlabel{section-trace@cref}{{[originalsection][0][]0FK6}{[1][0][]source}{}{}{}}
\makeatother
\crefname{originalsection}{original section}{original sections}
\Crefname{originalsection}{Original section}{Original sections}
\makeatletter
\newlabel{more-algebra-section-determinants-complexes}{{0FJI}{source}{}{sourcepointer-3}{}}
\newlabel{more-algebra-section-determinants-complexes@cref}{{[originalsection][0][]0FJI}{[1][0][]source}{}{}{}}
\makeatother
\crefname{originaldefinition}{original definition}{original definitions}
\Crefname{originaldefinition}{Original definition}{Original definitions}
\makeatletter
\newlabel{algebra-definition-relative-global-complete-intersection}{{00SP}{source}{}{sourcepointer-4}{}}
\newlabel{algebra-definition-relative-global-complete-intersection@cref}{{[originaldefinition][0][]00SP}{[1][0][]source}{}{}{}}
\makeatother
\crefname{originallemma}{original lemma}{original lemmas}
\Crefname{originallemma}{Original lemma}{Original lemmas}
\makeatletter
\newlabel{algebra-lemma-relative-global-complete-intersection-conormal}{{00SV}{source}{}{sourcepointer-5}{}}
\newlabel{algebra-lemma-relative-global-complete-intersection-conormal@cref}{{[originallemma][0][]00SV}{[1][0][]source}{}{}{}}
\makeatother
\crefname{originallemma}{original lemma}{original lemmas}
\Crefname{originallemma}{Original lemma}{Original lemmas}
\makeatletter
\newlabel{algebra-lemma-differential-seq}{{00RU}{source}{}{sourcepointer-6}{}}
\newlabel{algebra-lemma-differential-seq@cref}{{[originallemma][0][]00RU}{[1][0][]source}{}{}{}}
\makeatother
\crefname{originalsection}{original section}{original sections}
\Crefname{originalsection}{Original section}{Original sections}
\makeatletter
\newlabel{section-de-rham-complex}{{07HX}{source}{}{sourcepointer-7}{}}
\newlabel{section-de-rham-complex@cref}{{[originalsection][0][]07HX}{[1][0][]source}{}{}{}}
\makeatother
\crefname{originallemma}{original lemma}{original lemmas}
\Crefname{originallemma}{Original lemma}{Original lemmas}
\makeatletter
\newlabel{discriminant-lemma-syntomic-quasi-finite}{{0BWE}{source}{}{sourcepointer-8}{}}
\newlabel{discriminant-lemma-syntomic-quasi-finite@cref}{{[originallemma][0][]0BWE}{[1][0][]source}{}{}{}}
\makeatother
\crefname{originallemma}{original lemma}{original lemmas}
\Crefname{originallemma}{Original lemma}{Original lemmas}
\makeatletter
\newlabel{algebra-lemma-base-change-relative-global-complete-intersection}{{00SS}{source}{}{sourcepointer-9}{}}
\newlabel{algebra-lemma-base-change-relative-global-complete-intersection@cref}{{[originallemma][0][]00SS}{[1][0][]source}{}{}{}}
\makeatother
\crefname{originalproposition}{original proposition}{original propositions}
\Crefname{originalproposition}{Original proposition}{Original propositions}
\makeatletter
\newlabel{discriminant-proposition-tate-map}{{0FKD}{source}{}{sourcepointer-10}{}}
\newlabel{discriminant-proposition-tate-map@cref}{{[originalproposition][0][]0FKD}{[1][0][]source}{}{}{}}
\makeatother
\crefname{originalsection}{original section}{original sections}
\Crefname{originalsection}{Original section}{Original sections}
\makeatletter
\newlabel{discriminant-example-universal-quasi-finite-syntomic}{{0FK8}{source}{}{sourcepointer-11}{}}
\newlabel{discriminant-example-universal-quasi-finite-syntomic@cref}{{[originalsection][0][]0FK8}{[1][0][]source}{}{}{}}
\makeatother
\crefname{originallemma}{original lemma}{original lemmas}
\Crefname{originallemma}{Original lemma}{Original lemmas}
\makeatletter
\newlabel{discriminant-lemma-universal-quasi-finite-syntomic-etale}{{0FK9}{source}{}{sourcepointer-12}{}}
\newlabel{discriminant-lemma-universal-quasi-finite-syntomic-etale@cref}{{[originallemma][0][]0FK9}{[1][0][]source}{}{}{}}
\makeatother
\crefname{originalsection}{original section}{original sections}
\Crefname{originalsection}{Original section}{Original sections}
\makeatletter
\newlabel{discriminant-remark-local-description-delta}{{0FKC}{source}{}{sourcepointer-13}{}}
\newlabel{discriminant-remark-local-description-delta@cref}{{[originalsection][0][]0FKC}{[1][0][]source}{}{}{}}
\makeatother
\crefname{originallemma}{original lemma}{original lemmas}
\Crefname{originallemma}{Original lemma}{Original lemmas}
\makeatletter
\newlabel{morphisms-lemma-etale-at-point}{{02GU}{source}{}{sourcepointer-14}{}}
\newlabel{morphisms-lemma-etale-at-point@cref}{{[originallemma][0][]02GU}{[1][0][]source}{}{}{}}
\makeatother
\crefname{originallemma}{original lemma}{original lemmas}
\Crefname{originallemma}{Original lemma}{Original lemmas}
\makeatletter
\newlabel{discriminant-lemma-locally-comes-from-universal}{{0FKA}{source}{}{sourcepointer-15}{}}
\newlabel{discriminant-lemma-locally-comes-from-universal@cref}{{[originallemma][0][]0FKA}{[1][0][]source}{}{}{}}
\makeatother
\crefname{originallemma}{original lemma}{original lemmas}
\Crefname{originallemma}{Original lemma}{Original lemmas}
\makeatletter
\newlabel{properties-lemma-standard-open-two-affines}{{0H9B}{source}{}{sourcepointer-16}{}}
\newlabel{properties-lemma-standard-open-two-affines@cref}{{[originallemma][0][]0H9B}{[1][0][]source}{}{}{}}
\makeatother
\crefname{originallemma}{original lemma}{original lemmas}
\Crefname{originallemma}{Original lemma}{Original lemmas}
\makeatletter
\newlabel{more-algebra-lemma-module-over-fibre-product}{{07RU}{source}{}{sourcepointer-17}{}}
\newlabel{more-algebra-lemma-module-over-fibre-product@cref}{{[originallemma][0][]07RU}{[1][0][]source}{}{}{}}
\makeatother
\crefname{originallemma}{original lemma}{original lemmas}
\Crefname{originallemma}{Original lemma}{Original lemmas}
\makeatletter
\newlabel{more-algebra-lemma-properties-algebras-over-fibre-product}{{08KQ}{source}{}{sourcepointer-18}{}}
\newlabel{more-algebra-lemma-properties-algebras-over-fibre-product@cref}{{[originallemma][0][]08KQ}{[1][0][]source}{}{}{}}
\makeatother
\crefname{originallemma}{original lemma}{original lemmas}
\Crefname{originallemma}{Original lemma}{Original lemmas}
\makeatletter
\newlabel{more-morphisms-lemma-flat-lci}{{069K}{source}{}{sourcepointer-19}{}}
\newlabel{more-morphisms-lemma-flat-lci@cref}{{[originallemma][0][]069K}{[1][0][]source}{}{}{}}
\makeatother
\crefname{originallemma}{original lemma}{original lemmas}
\Crefname{originallemma}{Original lemma}{Original lemmas}
\makeatletter
\newlabel{morphisms-lemma-syntomic-flat-fibres}{{01UF}{source}{}{sourcepointer-20}{}}
\newlabel{morphisms-lemma-syntomic-flat-fibres@cref}{{[originallemma][0][]01UF}{[1][0][]source}{}{}{}}
\makeatother
\crefname{originallemma}{original lemma}{original lemmas}
\Crefname{originallemma}{Original lemma}{Original lemmas}
\makeatletter
\newlabel{algebra-lemma-relative-global-complete-intersection}{{00SW}{source}{}{sourcepointer-21}{}}
\newlabel{algebra-lemma-relative-global-complete-intersection@cref}{{[originallemma][0][]00SW}{[1][0][]source}{}{}{}}
\makeatother
\crefname{originallemma}{original lemma}{original lemmas}
\Crefname{originallemma}{Original lemma}{Original lemmas}
\makeatletter
\newlabel{algebra-lemma-isolated-point-fibre}{{00PK}{source}{}{sourcepointer-22}{}}
\newlabel{algebra-lemma-isolated-point-fibre@cref}{{[originallemma][0][]00PK}{[1][0][]source}{}{}{}}
\makeatother
\crefname{originallemma}{original lemma}{original lemmas}
\Crefname{originallemma}{Original lemma}{Original lemmas}
\makeatletter
\newlabel{algebra-lemma-syntomic}{{00SY}{source}{}{sourcepointer-23}{}}
\newlabel{algebra-lemma-syntomic@cref}{{[originallemma][0][]00SY}{[1][0][]source}{}{}{}}
\makeatother
\crefname{originallemma}{original lemma}{original lemmas}
\Crefname{originallemma}{Original lemma}{Original lemmas}
\makeatletter
\newlabel{morphisms-lemma-syntomic-locally-standard-syntomic}{{01UE}{source}{}{sourcepointer-24}{}}
\newlabel{morphisms-lemma-syntomic-locally-standard-syntomic@cref}{{[originallemma][0][]01UE}{[1][0][]source}{}{}{}}
\makeatother
\crefname{originallemma}{original lemma}{original lemmas}
\Crefname{originallemma}{Original lemma}{Original lemmas}
\makeatletter
\newlabel{more-morphisms-lemma-flat-fp-NL-lci}{{0FK3}{source}{}{sourcepointer-25}{}}
\newlabel{more-morphisms-lemma-flat-fp-NL-lci@cref}{{[originallemma][0][]0FK3}{[1][0][]source}{}{}{}}
\makeatother
\crefname{originallemma}{original lemma}{original lemmas}
\Crefname{originallemma}{Original lemma}{Original lemmas}
\makeatletter
\newlabel{lemma-dualizing-base-flat-flat}{{0BVE}{source}{}{sourcepointer-26}{}}
\newlabel{lemma-dualizing-base-flat-flat@cref}{{[originallemma][0][]0BVE}{[1][0][]source}{}{}{}}
\makeatother
\crefname{originallemma}{original lemma}{original lemmas}
\Crefname{originallemma}{Original lemma}{Original lemmas}
\makeatletter
\newlabel{more-morphisms-lemma-flat-and-free-at-point-fibre}{{080Q}{source}{}{sourcepointer-27}{}}
\newlabel{more-morphisms-lemma-flat-and-free-at-point-fibre@cref}{{[originallemma][0][]080Q}{[1][0][]source}{}{}{}}
\makeatother
\crefname{originallemma}{original lemma}{original lemmas}
\Crefname{originallemma}{Original lemma}{Original lemmas}
\makeatletter
\newlabel{lemma-dualizing-base-change-of-flat}{{0BVF}{source}{}{sourcepointer-28}{}}
\newlabel{lemma-dualizing-base-change-of-flat@cref}{{[originallemma][0][]0BVF}{[1][0][]source}{}{}{}}
\makeatother
\crefname{originalsection}{original section}{original sections}
\Crefname{originalsection}{Original section}{Original sections}
\makeatletter
\newlabel{remark-relative-dualizing-for-quasi-finite}{{0BVG}{source}{}{sourcepointer-29}{}}
\newlabel{remark-relative-dualizing-for-quasi-finite@cref}{{[originalsection][0][]0BVG}{[1][0][]source}{}{}{}}
\makeatother
\crefname{originallemma}{original lemma}{original lemmas}
\Crefname{originallemma}{Original lemma}{Original lemmas}
\makeatletter
\newlabel{algebra-lemma-Noether-normalization}{{00OY}{source}{}{sourcepointer-30}{}}
\newlabel{algebra-lemma-Noether-normalization@cref}{{[originallemma][0][]00OY}{[1][0][]source}{}{}{}}
\makeatother
\crefname{originallemma}{original lemma}{original lemmas}
\Crefname{originallemma}{Original lemma}{Original lemmas}
\makeatletter
\newlabel{dualizing-lemma-dualizing-finite}{{0AX0}{source}{}{sourcepointer-31}{}}
\newlabel{dualizing-lemma-dualizing-finite@cref}{{[originallemma][0][]0AX0}{[1][0][]source}{}{}{}}
\makeatother
\crefname{originallemma}{original lemma}{original lemmas}
\Crefname{originallemma}{Original lemma}{Original lemmas}
\makeatletter
\newlabel{dualizing-lemma-RHom-ext}{{0A71}{source}{}{sourcepointer-32}{}}
\newlabel{dualizing-lemma-RHom-ext@cref}{{[originallemma][0][]0A71}{[1][0][]source}{}{}{}}
\makeatother
\crefname{originallemma}{original lemma}{original lemmas}
\Crefname{originallemma}{Original lemma}{Original lemmas}
\makeatletter
\newlabel{dualizing-lemma-gorenstein}{{0DW9}{source}{}{sourcepointer-33}{}}
\newlabel{dualizing-lemma-gorenstein@cref}{{[originallemma][0][]0DW9}{[1][0][]source}{}{}{}}
\makeatother
\crefname{originallemma}{original lemma}{original lemmas}
\Crefname{originallemma}{Original lemma}{Original lemmas}
\makeatletter
\newlabel{dualizing-lemma-gorenstein-lci}{{0DWA}{source}{}{sourcepointer-34}{}}
\newlabel{dualizing-lemma-gorenstein-lci@cref}{{[originallemma][0][]0DWA}{[1][0][]source}{}{}{}}
\makeatother
\crefname{originallemma}{original lemma}{original lemmas}
\Crefname{originallemma}{Original lemma}{Original lemmas}
\makeatletter
\newlabel{algebra-lemma-quasi-finite-open}{{00QA}{source}{}{sourcepointer-35}{}}
\newlabel{algebra-lemma-quasi-finite-open@cref}{{[originallemma][0][]00QA}{[1][0][]source}{}{}{}}
\makeatother
\crefname{originallemma}{original lemma}{original lemmas}
\Crefname{originallemma}{Original lemma}{Original lemmas}
\makeatletter
\newlabel{perfect-lemma-determinant-two-term-complexes}{{0FJX}{source}{}{sourcepointer-36}{}}
\newlabel{perfect-lemma-determinant-two-term-complexes@cref}{{[originallemma][0][]0FJX}{[1][0][]source}{}{}{}}
\makeatother
\crefname{originallemma}{original lemma}{original lemmas}
\Crefname{originallemma}{Original lemma}{Original lemmas}
\makeatletter
\newlabel{more-morphisms-lemma-base-change-NL-flat}{{0FK0}{source}{}{sourcepointer-37}{}}
\newlabel{more-morphisms-lemma-base-change-NL-flat@cref}{{[originallemma][0][]0FK0}{[1][0][]source}{}{}{}}
\makeatother
\crefname{originalsection}{original section}{original sections}
\Crefname{originalsection}{Original section}{Original sections}
\makeatletter
\newlabel{perfect-remark-functorial-det}{{0FJY}{source}{}{sourcepointer-38}{}}
\newlabel{perfect-remark-functorial-det@cref}{{[originalsection][0][]0FJY}{[1][0][]source}{}{}{}}
\makeatother
\crefname{originalsection}{original section}{original sections}
\Crefname{originalsection}{Original section}{Original sections}
\makeatletter
\newlabel{remark-relative-dualizing-for-flat-quasi-finite}{{0BVJ}{source}{}{sourcepointer-39}{}}
\newlabel{remark-relative-dualizing-for-flat-quasi-finite@cref}{{[originalsection][0][]0BVJ}{[1][0][]source}{}{}{}}
\makeatother
\crefname{originalequation}{original equation}{original equations}
\Crefname{originalequation}{Original equation}{Original equations}
\makeatletter
\newlabel{equation-bc-dualizing}{{0BVB}{source}{}{sourcepointer-40}{}}
\newlabel{equation-bc-dualizing@cref}{{[originalequation][0][]0BVB}{[1][0][]source}{}{}{}}
\makeatother
\crefname{originallemma}{original lemma}{original lemmas}
\Crefname{originallemma}{Original lemma}{Original lemmas}
\makeatletter
\newlabel{lemma-trace-base-change}{{0BT9}{source}{}{sourcepointer-41}{}}
\newlabel{lemma-trace-base-change@cref}{{[originallemma][0][]0BT9}{[1][0][]source}{}{}{}}
\makeatother
\title{Canonical Garel maps on differential forms for quasi-finite syntomic morphisms}
\author{The Stacks Project Authors}
\date{stacks-project:derham.tex:local-snapshot-20261011; GNU FDL 1.2 or later}
\begin{document}
\maketitle
\noindent\textit{Editorial note.} The old dump stopped at3741; the exact cached chapter includes3742–3764 and the final universal-refinement step. The task uses the complete original range.\par\medskip
\noindent\textit{Editorial note.} The closing diagram uses the already proved local-universal-family construction; it is an author-compressed descent step, not an inserted alternative proof.\par\medskip
\noindent\textit{Editorial note.} The separate downstream trace map lemma-Garel is not counted here. The earlier Tate map is only a comparison of proof strategy; no free Garel map is imported.\par\medskip
\noindent\textit{Allowed standard prerequisites (editorial).}\begin{itemize}
\item Explicit prerequisite boundary. Differentials, exterior powers and determinant lines; syntomic local complete-intersection presentations, etale locus, fibre-product algebra patching, and standard affine localization. The local funny-map construction, compatibility lemmas and universal family are retained author obligations.
\end{itemize}\medskip
\ifdefined\hypertarget\hypertarget{sourcepointer-1}{}\fi
\noindent\textit{Source reference.} Original section 0FKB: focused-seeds/AGP/sources/stacks-snapshot-20261011/tags.txt, source lines 19456--19456. Stable Stacks tag for an expressly allowed foundational or notational input; the counted construction remains the two concrete core obligations. The retained local core is not replaced by this pointer. See standard\_prerequisites.\par\medskip
\ifdefined\hypertarget\hypertarget{sourcepointer-2}{}\fi
\noindent\textit{Source reference.} Original section 0FK6: focused-seeds/AGP/sources/stacks-snapshot-20261011/tags.txt, source lines 19451--19451. Stable Stacks tag for an expressly allowed foundational or notational input; the counted construction remains the two concrete core obligations. The retained local core is not replaced by this pointer. See standard\_prerequisites.\par\medskip
\ifdefined\hypertarget\hypertarget{sourcepointer-3}{}\fi
\noindent\textit{Source reference.} Original section 0FJI: focused-seeds/AGP/sources/stacks-snapshot-20261011/tags.txt, source lines 19428--19428. Stable Stacks tag for an expressly allowed foundational or notational input; the counted construction remains the two concrete core obligations. The retained local core is not replaced by this pointer. See standard\_prerequisites.\par\medskip
\ifdefined\hypertarget\hypertarget{sourcepointer-4}{}\fi
\noindent\textit{Source reference.} Original definition 00SP: focused-seeds/AGP/sources/stacks-snapshot-20261011/tags.txt, source lines 1047--1047. Stable Stacks tag for an expressly allowed foundational or notational input; the counted construction remains the two concrete core obligations. The retained local core is not replaced by this pointer. See standard\_prerequisites.\par\medskip
\ifdefined\hypertarget\hypertarget{sourcepointer-5}{}\fi
\noindent\textit{Source reference.} Original lemma 00SV: focused-seeds/AGP/sources/stacks-snapshot-20261011/tags.txt, source lines 1053--1053. Stable Stacks tag for an expressly allowed foundational or notational input; the counted construction remains the two concrete core obligations. The retained local core is not replaced by this pointer. See standard\_prerequisites.\par\medskip
\ifdefined\hypertarget\hypertarget{sourcepointer-6}{}\fi
\noindent\textit{Source reference.} Original lemma 00RU: focused-seeds/AGP/sources/stacks-snapshot-20261011/tags.txt, source lines 1016--1016. Stable Stacks tag for an expressly allowed foundational or notational input; the counted construction remains the two concrete core obligations. The retained local core is not replaced by this pointer. See standard\_prerequisites.\par\medskip
\ifdefined\hypertarget\hypertarget{sourcepointer-7}{}\fi
\noindent\textit{Source reference.} Original section 07HX: focused-seeds/AGP/sources/stacks-snapshot-20261011/tags.txt, source lines 9567--9567. Stable Stacks tag for an expressly allowed foundational or notational input; the counted construction remains the two concrete core obligations. The retained local core is not replaced by this pointer. See standard\_prerequisites.\par\medskip
\ifdefined\hypertarget\hypertarget{sourcepointer-8}{}\fi
\noindent\textit{Source reference.} Original lemma 0BWE: focused-seeds/AGP/sources/stacks-snapshot-20261011/tags.txt, source lines 14942--14942. Stable Stacks tag for an expressly allowed foundational or notational input; the counted construction remains the two concrete core obligations. The retained local core is not replaced by this pointer. See standard\_prerequisites.\par\medskip
\ifdefined\hypertarget\hypertarget{sourcepointer-9}{}\fi
\noindent\textit{Source reference.} Original lemma 00SS: focused-seeds/AGP/sources/stacks-snapshot-20261011/tags.txt, source lines 1050--1050. Stable Stacks tag for an expressly allowed foundational or notational input; the counted construction remains the two concrete core obligations. The retained local core is not replaced by this pointer. See standard\_prerequisites.\par\medskip
\ifdefined\hypertarget\hypertarget{sourcepointer-10}{}\fi
\noindent\textit{Source reference.} Original proposition 0FKD: focused-seeds/AGP/sources/stacks-snapshot-20261011/tags.txt, source lines 19458--19458. Stable Stacks tag for an expressly allowed foundational or notational input; the counted construction remains the two concrete core obligations. The retained local core is not replaced by this pointer. See standard\_prerequisites.\par\medskip
\ifdefined\hypertarget\hypertarget{sourcepointer-11}{}\fi
\noindent\textit{Source reference.} Original section 0FK8: focused-seeds/AGP/sources/stacks-snapshot-20261011/tags.txt, source lines 19453--19453. Stable Stacks tag for an expressly allowed foundational or notational input; the counted construction remains the two concrete core obligations. The retained local core is not replaced by this pointer. See standard\_prerequisites.\par\medskip
\ifdefined\hypertarget\hypertarget{sourcepointer-12}{}\fi
\noindent\textit{Source reference.} Original lemma 0FK9: focused-seeds/AGP/sources/stacks-snapshot-20261011/tags.txt, source lines 19454--19454. Stable Stacks tag for an expressly allowed foundational or notational input; the counted construction remains the two concrete core obligations. The retained local core is not replaced by this pointer. See standard\_prerequisites.\par\medskip
\ifdefined\hypertarget\hypertarget{sourcepointer-13}{}\fi
\noindent\textit{Source reference.} Original section 0FKC: focused-seeds/AGP/sources/stacks-snapshot-20261011/tags.txt, source lines 19457--19457. Stable Stacks tag for an expressly allowed foundational or notational input; the counted construction remains the two concrete core obligations. The retained local core is not replaced by this pointer. See standard\_prerequisites.\par\medskip
\ifdefined\hypertarget\hypertarget{sourcepointer-14}{}\fi
\noindent\textit{Source reference.} Original lemma 02GU: focused-seeds/AGP/sources/stacks-snapshot-20261011/tags.txt, source lines 3241--3241. Stable Stacks tag for an expressly allowed foundational or notational input; the counted construction remains the two concrete core obligations. The retained local core is not replaced by this pointer. See standard\_prerequisites.\par\medskip
\ifdefined\hypertarget\hypertarget{sourcepointer-15}{}\fi
\noindent\textit{Source reference.} Original lemma 0FKA: focused-seeds/AGP/sources/stacks-snapshot-20261011/tags.txt, source lines 19455--19455. Stable Stacks tag for an expressly allowed foundational or notational input; the counted construction remains the two concrete core obligations. The retained local core is not replaced by this pointer. See standard\_prerequisites.\par\medskip
\ifdefined\hypertarget\hypertarget{sourcepointer-16}{}\fi
\noindent\textit{Source reference.} Original lemma 0H9B: focused-seeds/AGP/sources/stacks-snapshot-20261011/tags.txt, source lines 21521--21521. Stable Stacks tag for an expressly allowed foundational or notational input; the counted construction remains the two concrete core obligations. The retained local core is not replaced by this pointer. See standard\_prerequisites.\par\medskip
\ifdefined\hypertarget\hypertarget{sourcepointer-17}{}\fi
\noindent\textit{Source reference.} Original lemma 07RU: focused-seeds/AGP/sources/stacks-snapshot-20261011/tags.txt, source lines 9879--9879. Stable Stacks tag for an expressly allowed foundational or notational input; the counted construction remains the two concrete core obligations. The retained local core is not replaced by this pointer. See standard\_prerequisites.\par\medskip
\ifdefined\hypertarget\hypertarget{sourcepointer-18}{}\fi
\noindent\textit{Source reference.} Original lemma 08KQ: focused-seeds/AGP/sources/stacks-snapshot-20261011/tags.txt, source lines 10893--10893. Stable Stacks tag for an expressly allowed foundational or notational input; the counted construction remains the two concrete core obligations. The retained local core is not replaced by this pointer. See standard\_prerequisites.\par\medskip
\ifdefined\hypertarget\hypertarget{sourcepointer-19}{}\fi
\noindent\textit{Source reference.} Original lemma 069K: focused-seeds/AGP/sources/stacks-snapshot-20261011/tags.txt, source lines 8044--8044. Stable Stacks tag for an expressly allowed foundational or notational input; the counted construction remains the two concrete core obligations. The retained local core is not replaced by this pointer. See standard\_prerequisites.\par\medskip
\ifdefined\hypertarget\hypertarget{sourcepointer-20}{}\fi
\noindent\textit{Source reference.} Original lemma 01UF: focused-seeds/AGP/sources/stacks-snapshot-20261011/tags.txt, source lines 2408--2408. Stable Stacks tag for an expressly allowed foundational or notational input; the counted construction remains the two concrete core obligations. The retained local core is not replaced by this pointer. See standard\_prerequisites.\par\medskip
\ifdefined\hypertarget\hypertarget{sourcepointer-21}{}\fi
\noindent\textit{Source reference.} Original lemma 00SW: focused-seeds/AGP/sources/stacks-snapshot-20261011/tags.txt, source lines 1054--1054. Stable Stacks tag for an expressly allowed foundational or notational input; the counted construction remains the two concrete core obligations. The retained local core is not replaced by this pointer. See standard\_prerequisites.\par\medskip
\ifdefined\hypertarget\hypertarget{sourcepointer-22}{}\fi
\noindent\textit{Source reference.} Original lemma 00PK: focused-seeds/AGP/sources/stacks-snapshot-20261011/tags.txt, source lines 934--934. Stable Stacks tag for an expressly allowed foundational or notational input; the counted construction remains the two concrete core obligations. The retained local core is not replaced by this pointer. See standard\_prerequisites.\par\medskip
\ifdefined\hypertarget\hypertarget{sourcepointer-23}{}\fi
\noindent\textit{Source reference.} Original lemma 00SY: focused-seeds/AGP/sources/stacks-snapshot-20261011/tags.txt, source lines 1056--1056. Stable Stacks tag for an expressly allowed foundational or notational input; the counted construction remains the two concrete core obligations. The retained local core is not replaced by this pointer. See standard\_prerequisites.\par\medskip
\ifdefined\hypertarget\hypertarget{sourcepointer-24}{}\fi
\noindent\textit{Source reference.} Original lemma 01UE: focused-seeds/AGP/sources/stacks-snapshot-20261011/tags.txt, source lines 2407--2407. Stable Stacks tag for an expressly allowed foundational or notational input; the counted construction remains the two concrete core obligations. The retained local core is not replaced by this pointer. See standard\_prerequisites.\par\medskip
\ifdefined\hypertarget\hypertarget{sourcepointer-25}{}\fi
\noindent\textit{Source reference.} Original lemma 0FK3: focused-seeds/AGP/sources/stacks-snapshot-20261011/tags.txt, source lines 19448--19448. Stable Stacks tag for an expressly allowed foundational or notational input; the counted construction remains the two concrete core obligations. The retained local core is not replaced by this pointer. See standard\_prerequisites.\par\medskip
\ifdefined\hypertarget\hypertarget{sourcepointer-26}{}\fi
\noindent\textit{Source reference.} Original lemma 0BVE: focused-seeds/AGP/sources/stacks-snapshot-20261011/tags.txt, source lines 14907--14907. Stable Stacks tag for an expressly allowed foundational or notational input; the counted construction remains the two concrete core obligations. The retained local core is not replaced by this pointer. See standard\_prerequisites.\par\medskip
\ifdefined\hypertarget\hypertarget{sourcepointer-27}{}\fi
\noindent\textit{Source reference.} Original lemma 080Q: focused-seeds/AGP/sources/stacks-snapshot-20261011/tags.txt, source lines 10190--10190. Stable Stacks tag for an expressly allowed foundational or notational input; the counted construction remains the two concrete core obligations. The retained local core is not replaced by this pointer. See standard\_prerequisites.\par\medskip
\ifdefined\hypertarget\hypertarget{sourcepointer-28}{}\fi
\noindent\textit{Source reference.} Original lemma 0BVF: focused-seeds/AGP/sources/stacks-snapshot-20261011/tags.txt, source lines 14908--14908. Stable Stacks tag for an expressly allowed foundational or notational input; the counted construction remains the two concrete core obligations. The retained local core is not replaced by this pointer. See standard\_prerequisites.\par\medskip
\ifdefined\hypertarget\hypertarget{sourcepointer-29}{}\fi
\noindent\textit{Source reference.} Original section 0BVG: focused-seeds/AGP/sources/stacks-snapshot-20261011/tags.txt, source lines 14909--14909. Stable Stacks tag for an expressly allowed foundational or notational input; the counted construction remains the two concrete core obligations. The retained local core is not replaced by this pointer. See standard\_prerequisites.\par\medskip
\ifdefined\hypertarget\hypertarget{sourcepointer-30}{}\fi
\noindent\textit{Source reference.} Original lemma 00OY: focused-seeds/AGP/sources/stacks-snapshot-20261011/tags.txt, source lines 912--912. Stable Stacks tag for an expressly allowed foundational or notational input; the counted construction remains the two concrete core obligations. The retained local core is not replaced by this pointer. See standard\_prerequisites.\par\medskip
\ifdefined\hypertarget\hypertarget{sourcepointer-31}{}\fi
\noindent\textit{Source reference.} Original lemma 0AX0: focused-seeds/AGP/sources/stacks-snapshot-20261011/tags.txt, source lines 13738--13738. Stable Stacks tag for an expressly allowed foundational or notational input; the counted construction remains the two concrete core obligations. The retained local core is not replaced by this pointer. See standard\_prerequisites.\par\medskip
\ifdefined\hypertarget\hypertarget{sourcepointer-32}{}\fi
\noindent\textit{Source reference.} Original lemma 0A71: focused-seeds/AGP/sources/stacks-snapshot-20261011/tags.txt, source lines 12864--12864. Stable Stacks tag for an expressly allowed foundational or notational input; the counted construction remains the two concrete core obligations. The retained local core is not replaced by this pointer. See standard\_prerequisites.\par\medskip
\ifdefined\hypertarget\hypertarget{sourcepointer-33}{}\fi
\noindent\textit{Source reference.} Original lemma 0DW9: focused-seeds/AGP/sources/stacks-snapshot-20261011/tags.txt, source lines 17389--17389. Stable Stacks tag for an expressly allowed foundational or notational input; the counted construction remains the two concrete core obligations. The retained local core is not replaced by this pointer. See standard\_prerequisites.\par\medskip
\ifdefined\hypertarget\hypertarget{sourcepointer-34}{}\fi
\noindent\textit{Source reference.} Original lemma 0DWA: focused-seeds/AGP/sources/stacks-snapshot-20261011/tags.txt, source lines 17390--17390. Stable Stacks tag for an expressly allowed foundational or notational input; the counted construction remains the two concrete core obligations. The retained local core is not replaced by this pointer. See standard\_prerequisites.\par\medskip
\ifdefined\hypertarget\hypertarget{sourcepointer-35}{}\fi
\noindent\textit{Source reference.} Original lemma 00QA: focused-seeds/AGP/sources/stacks-snapshot-20261011/tags.txt, source lines 960--960. Stable Stacks tag for an expressly allowed foundational or notational input; the counted construction remains the two concrete core obligations. The retained local core is not replaced by this pointer. See standard\_prerequisites.\par\medskip
\ifdefined\hypertarget\hypertarget{sourcepointer-36}{}\fi
\noindent\textit{Source reference.} Original lemma 0FJX: focused-seeds/AGP/sources/stacks-snapshot-20261011/tags.txt, source lines 19442--19442. Stable Stacks tag for an expressly allowed foundational or notational input; the counted construction remains the two concrete core obligations. The retained local core is not replaced by this pointer. See standard\_prerequisites.\par\medskip
\ifdefined\hypertarget\hypertarget{sourcepointer-37}{}\fi
\noindent\textit{Source reference.} Original lemma 0FK0: focused-seeds/AGP/sources/stacks-snapshot-20261011/tags.txt, source lines 19445--19445. Stable Stacks tag for an expressly allowed foundational or notational input; the counted construction remains the two concrete core obligations. The retained local core is not replaced by this pointer. See standard\_prerequisites.\par\medskip
\ifdefined\hypertarget\hypertarget{sourcepointer-38}{}\fi
\noindent\textit{Source reference.} Original section 0FJY: focused-seeds/AGP/sources/stacks-snapshot-20261011/tags.txt, source lines 19443--19443. Stable Stacks tag for an expressly allowed foundational or notational input; the counted construction remains the two concrete core obligations. The retained local core is not replaced by this pointer. See standard\_prerequisites.\par\medskip
\ifdefined\hypertarget\hypertarget{sourcepointer-39}{}\fi
\noindent\textit{Source reference.} Original section 0BVJ: focused-seeds/AGP/sources/stacks-snapshot-20261011/tags.txt, source lines 14912--14912. Stable Stacks tag for an expressly allowed foundational or notational input; the counted construction remains the two concrete core obligations. The retained local core is not replaced by this pointer. See standard\_prerequisites.\par\medskip
\ifdefined\hypertarget\hypertarget{sourcepointer-40}{}\fi
\noindent\textit{Source reference.} Original equation 0BVB: focused-seeds/AGP/sources/stacks-snapshot-20261011/tags.txt, source lines 14904--14904. Stable Stacks tag for an expressly allowed foundational or notational input; the counted construction remains the two concrete core obligations. The retained local core is not replaced by this pointer. See standard\_prerequisites.\par\medskip
\ifdefined\hypertarget\hypertarget{sourcepointer-41}{}\fi
\noindent\textit{Source reference.} Original lemma 0BT9: focused-seeds/AGP/sources/stacks-snapshot-20261011/tags.txt, source lines 14832--14832. Stable Stacks tag for an expressly allowed foundational or notational input; the counted construction remains the two concrete core obligations. The retained local core is not replaced by this pointer. See standard\_prerequisites.\par\medskip
\begin{lemma}
\label{lemma-Garel-upstairs}
There exists a unique rule that to every locally quasi-finite syntomic
morphism of schemes $f : Y \to X$ assigns $\mathcal{O}_Y$-module maps
$$
c^p_{Y/X} :
\Omega^p_{Y/\mathbf{Z}}
\longrightarrow
f^*\Omega^p_{X/\mathbf{Z}} \otimes_{\mathcal{O}_Y} \det(\NL_{Y/X})
$$
for $p \geq 0$ satisfying (\ref{item-degree-zero}), (\ref{item-multiplicative}),
and (\ref{item-base-change}). In particular, the composition of
$c^p_{Y/X}$ with $f^*\Omega^p_{X/\mathbf{Z}} \to \Omega^p_{Y/\mathbf{Z}}$
is multiplication by $\delta(\NL_{Y/X})$.
\end{lemma}

\section{The Tate map}
\label{section-tate-map}

\noindent
In this section we produce an isomorphism between
the determinant of the relative cotangent complex and
the relative dualizing module for a locally quasi-finite
syntomic morphism of locally Noetherian schemes. Following
\cite[1.4.4]{Garel} we dub the isomorphism the Tate map.
Our approach is to avoid doing local calculations as
much as is possible.

\medskip\noindent
Let $Y \to X$ be a locally quasi-finite syntomic morphism of schemes.
We will use all the equivalent conditions for this notion given in
Lemma \ref{lemma-syntomic-quasi-finite} without further mention in
this section. In particular, we see that $\NL_{Y/X}$ is a perfect
object of $D(\mathcal{O}_Y)$ with tor-amplitude in $[-1, 0]$. Thus
we have a canonical invertible module
$\det(\NL_{Y/X})$ on $Y$ and a global section
$$
\delta(\NL_{Y/X}) \in \Gamma(Y, \det(\NL_{Y/X}))
$$
See Derived Categories of Schemes, Lemma
\ref{perfect-lemma-determinant-two-term-complexes}.
Suppose given a commutative diagram of schemes
$$
\xymatrix{
Y' \ar[r]_b \ar[d] & Y \ar[d] \\
X' \ar[r] & X
}
$$
whose vertical arrows are locally quasi-finite syntomic and which
induces an isomorphism of $Y'$ with an open of $X' \times_X Y$.
Then the canonical map
$$
Lb^*\NL_{Y/X} \longrightarrow \NL_{Y'/X'}
$$
is a quasi-isomorphism by
More on Morphisms, Lemma \ref{more-morphisms-lemma-base-change-NL-flat}.
Thus we get a canonical isomorphism
$b^*\det(\NL_{Y/X}) \to \det(\NL_{Y'/X'})$ which sends the
canonical section $\delta(\NL_{Y/X})$ to $\delta(\NL_{Y'/ X'})$, see
Derived Categories of Schemes, Remark \ref{perfect-remark-functorial-det}.

\begin{remark}
\label{remark-local-description-delta}
Let $Y \to X$ be a locally quasi-finite syntomic morphism of schemes.
What does the pair $(\det(\NL_{Y/X}), \delta(\NL_{Y/X}))$ look
like locally? Choose affine opens $V = \Spec(B) \subset Y$,
$U = \Spec(A) \subset X$ with $f(V) \subset U$ and an integer $n$ and
$f_1, \ldots, f_n \in A[x_1, \ldots, x_n]$ such that
$B = A[x_1, \ldots, x_n]/(f_1, \ldots, f_n)$. Then
$$
\NL_{B/A} = \left(
(f_1, \ldots, f_n)/(f_1, \ldots, f_n)^2
\longrightarrow
\bigoplus\nolimits_{i = 1, \ldots, n} B \text{d} x_i\right)
$$
and $(f_1, \ldots, f_n)/(f_1, \ldots, f_n)^2$ is free with generators
the classes $\overline{f}_i$. See proof of
Lemma \ref{lemma-syntomic-quasi-finite}.
Thus $\det(L_{B/A})$ is free on the generator
$$
\text{d}x_1 \wedge \ldots \wedge \text{d}x_n
\otimes
(\overline{f}_1 \wedge \ldots \wedge \overline{f}_n)^{\otimes -1}
$$
and the section $\delta(\NL_{B/A})$ is the element
$$
\delta(\NL_{B/A}) =
\det(\partial f_j/ \partial x_i) \cdot
\text{d}x_1 \wedge \ldots \wedge \text{d}x_n
\otimes
(\overline{f}_1 \wedge \ldots \wedge \overline{f}_n)^{\otimes -1}
$$
by definition.
\end{remark}

\noindent
Let $Y \to X$ be a locally quasi-finite syntomic morphism of
locally Noetherian schemes. By
Remarks \ref{remark-relative-dualizing-for-quasi-finite} and
\ref{remark-relative-dualizing-for-flat-quasi-finite} we have
a coherent $\mathcal{O}_Y$-module $\omega_{Y/X}$ and a canonical
global section
$$
\tau_{Y/X} \in \Gamma(Y, \omega_{Y/X})
$$
which affine locally recovers the pair $\omega_{B/A}, \tau_{B/A}$.
By Lemma \ref{lemma-dualizing-syntomic-quasi-finite} the module
$\omega_{Y/X}$ is invertible. Suppose given a commutative diagram of
locally Noetherian schemes
$$
\xymatrix{
Y' \ar[r]_b \ar[d] & Y \ar[d] \\
X' \ar[r] & X
}
$$
whose vertical arrows are locally quasi-finite syntomic and which
induces an isomorphism of $Y'$ with an open of $X' \times_X Y$.
Then there is a canonical base change map
$$
b^*\omega_{Y/X} \longrightarrow \omega_{Y'/X'}
$$
which is an isomorphism
mapping $\tau_{Y/X}$ to $\tau_{Y'/X'}$. Namely, the base change map
in the affine setting is (\ref{equation-bc-dualizing}), it is an
isomorphism by Lemma \ref{lemma-dualizing-base-change-of-flat}, and it
maps $\tau_{Y/X}$ to $\tau_{Y'/X'}$ by
Lemma \ref{lemma-trace-base-change} part (1).

\section{Comparing sheaves of differential forms}
\label{section-quasi-finite-syntomic}

\noindent
The goal of this section is to construct for every locally quasi-finite
syntomic morphism of schemes $f : Y \to X$ maps
$$
c^p_{Y/X} :
\Omega^p_{Y/\mathbf{Z}}
\longrightarrow
f^*\Omega^p_{X/\mathbf{Z}} \otimes_{\mathcal{O}_Y} \det(\NL_{Y/X})
$$
for all $p \geq 0$ satisfying the following properties
\begin{enumerate}
\item
\label{item-degree-zero}
$c_{Y/X}^0(1) = \delta(\NL_{Y/X})$ see
Discriminants, Section \ref{discriminant-section-tate-map},
\item
\label{item-multiplicative}
$c_{Y/X}^{q + p}(\omega \wedge \eta) = \omega \wedge c_{Y/X}^p(\eta)$
for local sections $\omega$ of $f^*\Omega^q_{X/\mathbf{Z}}$
and $\eta$ of $\Omega^p_{Y/\mathbf{Z}}$,
\item
\label{item-base-change}
the formation of $c^p_{Y/X}$ is compatible with restriction to opens
and with base change
(as formulated in Lemma \ref{lemma-base-change-Garel-upstairs}).
\end{enumerate}
Note that these conditions imply that the maps $c^p_{Y/X}$ are
$\mathcal{O}_Y$-linear and that the composition with
$f^*\Omega^p_{X/\mathbf{Z}} \to \Omega^p_{Y/\mathbf{Z}}$
is multiplication by $\delta(\NL_{Y/X})$.
We will in fact show that there exists a unique collection of operators
$c^p_{Y/X}$ with properties (1), (2), (3). This result will be applied in
Section \ref{section-trace} to the construction of the trace map on
the de Rham complexes if $f$ is finite.

\begin{lemma}
\label{lemma-funny-map}
Let $R$ be a ring and consider a commutative diagram
$$
\xymatrix{
0 \ar[r] &
K^0 \ar[r] &
L^0 \ar[r] &
M^0 \ar[r] & 0 \\
& & L^{-1} \ar[u]_\partial \ar@{=}[r] &
M^{-1} \ar[u]
}
$$
of $R$-modules with exact top row and $M^0$ and $M^{-1}$
finite free of the same rank. For $p \geq 0$ there are canonical maps
$$
c^p : 
\wedge^p(H^0(L^\bullet))
\longrightarrow
\wedge^p(K^0) \otimes_R \det(M^\bullet)
$$
with the following properties:
\begin{enumerate}
\item $c^0(1) = \delta(M^\bullet)$, and
\item $c^{q + p}(\omega \wedge \eta) = \omega \wedge c^p(\eta)$
for $\omega \in \wedge^q(K^0)$ and $\eta \in \wedge^p(H^0(L^\bullet))$.
\end{enumerate}
In particular, the composition of $c^p$ with the map
$\wedge^p(K^0) \to \wedge^p(H^0(L^\bullet))$ is multiplication
with $\delta(M^\bullet)$.
\end{lemma}

\begin{proof}
Say $M^0$ and $M^{-1}$ are free of rank $n$. For every $p \geq 0$
there is a unique surjection
$$
\pi_p :
\wedge^{p + n}(L^0)
\longrightarrow
\wedge^p(K^0) \otimes \wedge^n(M^0)
$$
with the following two properties; (i) we have
$
\pi_p(k_1 \wedge \ldots \wedge k_p \wedge l_1 \wedge \ldots \wedge l_n) =
k_1 \wedge \ldots \wedge k_p \otimes
m_1 \wedge \ldots \wedge m_n
$
if $k_1, \ldots, k_p \in K^0$ and $l_1, \ldots, l_n \in L^0$ map
to $m_1, \ldots, m_n \in M^0$ and (ii) the kernel of $\pi_p$ is the
submodule generated by wedges
$l_1 \wedge \ldots \wedge l_{p + n}$ such that $> p$ of the
$l_j$ are in $K^0$. Pick a basis $e_1, \ldots, e_n$ for $L^{-1} = M^{-1}$.
Then we define our map by the rule
$$
\eta \mapsto
\pi_p(\tilde \eta \wedge \partial(e_1) \wedge \ldots \wedge \partial(e_n))
\otimes (e_1 \wedge \ldots \wedge e_n)^{\otimes -1}
$$
where $\eta \in \wedge^p(H^0(L^\bullet))$ and $\tilde \eta \in \wedge^p(L^0)$
is a lift of $\eta$. The expression on the right hand side makes sense because
$\det(M^\bullet) = \wedge^n(M^0) \otimes \wedge^n(M^{-1})^{\otimes -1}$.
The expression on the right is independent of the choice of $\tilde \eta$ since
$\partial(\xi) \wedge \partial(e_1) \wedge \ldots \wedge \partial(e_n) = 0$
for any $\xi \in L^{-1}$. A straightforward calculation shows that the map
is independent of the choice of the basis of $L^{-1} = M^{-1}$.
It is immediate from the definition of $\delta(M^\bullet)$ in
More on Algebra, Section \ref{more-algebra-section-determinants-complexes}
that $c^0(1) = \delta(M^\bullet)$. Finally, the rule in statement (2)
can be directly verified from the definition.
\end{proof}

\begin{lemma}
\label{lemma-funny-map-functorial}
Let $R_1 \to R_2$ be a ring homomorphism. For $i = 1, 2$ consider
commutative diagrams
$$
\xymatrix{
0 \ar[r] &
K_i^0 \ar[r] &
L_i^0 \ar[r] &
M_i^0 \ar[r] & 0 \\
& & L_i^{-1} \ar[u]_{\partial_i} \ar@{=}[r] &
M_i^{-1} \ar[u]
}
$$
of $R_i$-modules as in Lemma \ref{lemma-funny-map}. Assume we have
maps $K_1^0 \to K_2^0$, $L_1^j \to L_2^j$, $M_1^j \to M_2^j$
compatible with the given ring map $R_1 \to R_2$ and
compatible with the maps in the displayed diagrams. If the maps
$M_1^j \otimes_{R_1} R_2 \to M_2^j$ are isomorphisms, then
the diagrams
$$
\xymatrix{ 
\wedge^p(H^0(L_1^\bullet)) \ar[d]
\ar[r]_-{c^p} &
\wedge^p(K_1^0) \otimes_{R_1} \det(M_1^\bullet) \ar[d] \\
\wedge^p(H^0(L_2^\bullet))
\ar[r]^-{c^p} &
\wedge^p(K_2^0) \otimes_{R_2} \det(M_2^\bullet)
}
$$
commute.
\end{lemma}

\begin{proof}
This follows from the explicit description of the map given
in the proof of Lemma \ref{lemma-funny-map}. Note that we need
the assumption that $M_1^j \otimes_{R_1} R_2 \to M_2^j$ are isomorphisms
to see that the basis we use for $M_1^{-1}$ maps to a basis
for $M_2^{-1}$ (and hence in checking the compatibility in
both constructions we can use the ``same'' basis).
\end{proof}

\begin{remark}
\label{remark-local-description}
Let $A$ be a ring. Let $P = A[x_1, \ldots, x_n]$. Let
$f_1, \ldots, f_n \in P$ and set $B = P/(f_1, \ldots, f_n)$.
Assume $A \to B$ is quasi-finite. Then
$B$ is a relative global complete intersection over $A$ (Algebra, Definition
\ref{algebra-definition-relative-global-complete-intersection}) and
$(f_1, \ldots, f_n)/(f_1, \ldots, f_n)^2$ is free with generators
the classes $\overline{f}_i$ by Algebra, Lemma
\ref{algebra-lemma-relative-global-complete-intersection-conormal}.
Consider the following diagram
$$
\xymatrix{
\Omega_{A/\mathbf{Z}} \otimes_A B \ar[r] &
\Omega_{P/\mathbf{Z}} \otimes_P B \ar[r] &
\Omega_{P/A} \otimes_P B \\
&
(f_1, \ldots, f_n)/(f_1, \ldots, f_n)^2 \ar[u] \ar@{=}[r] &
(f_1, \ldots, f_n)/(f_1, \ldots, f_n)^2 \ar[u]
}
$$
The right column represents $\NL_{B/A}$ in $D(B)$ hence has cohomology
$\Omega_{B/A}$ in degree $0$. The top row is the split short exact sequence
$0 \to \Omega_{A/\mathbf{Z}} \otimes_A B \to
\Omega_{P/\mathbf{Z}} \otimes_P B \to \Omega_{P/A} \otimes_P B \to 0$.
The middle column has cohomology $\Omega_{B/\mathbf{Z}}$ in degree $0$
by Algebra, Lemma \ref{algebra-lemma-differential-seq}.
Thus by Lemma \ref{lemma-funny-map} we obtain canonical $B$-module maps
$$
\Omega^p_{B/\mathbf{Z}} \longrightarrow
\Omega^p_{A/\mathbf{Z}} \otimes_A \det(\NL_{B/A})
$$
for $p \geq 0$ satisfying properties (1) and (2) and in particular
satisfy the property that the composition with
$\Omega^p_{A/\mathbf{Z}} \to \Omega^p_{B/\mathbf{Z}}$
is multiplication by $\delta(\NL_{B/A})$.
\end{remark}

\begin{lemma}
\label{lemma-base-change-Garel-upstairs}
Consider a commutative diagram
$$
\xymatrix{
Y' \ar[d]_{f'} \ar[r]_b & Y \ar[d]^f \\
X' \ar[r]^a & X
}
$$
of schemes which induces an isomorphism of $Y'$ with an open subscheme of
$X' \times_X Y$. Assume $f$ is locally quasi-finite and syntomic and
assume given maps $c^p_{Y/X} : \Omega^p_{Y/\mathbf{Z}} \to
f^*\Omega^p_{X/\mathbf{Z}} \otimes_{\mathcal{O}_Y} \det(\NL_{Y/X})$
for $p \geq 0$ satisfying
(\ref{item-degree-zero}) and (\ref{item-multiplicative}).
Then there is at most one collection of maps
$c^p_{Y'/X'} : \Omega^p_{Y'/\mathbf{Z}} \to
(f')^*\Omega^p_{X'/\mathbf{Z}} \otimes_{\mathcal{O}_{Y'}} \det(\NL_{Y'/X'})$
for $p \geq 0$
satisfying (\ref{item-degree-zero}) and (\ref{item-multiplicative})
such that the diagrams
$$
\xymatrix{
b^*\Omega^p_{Y/\mathbf{Z}} \ar[rr]_-{b^*c^p_{Y/X}} \ar[d] & &
b^*(f^*\Omega^p_{X/\mathbf{Z}} \otimes_{\mathcal{O}_Y} \det(\NL_{Y/X}))
\ar@{=}[r] &
(f')^*a^*\Omega^p_{X/\mathbf{Z}} \otimes_{\mathcal{O}_Y} b^*\det(\NL_{Y/X}))
\ar[d] \\
\Omega^p_{Y'/\mathbf{Z}} \ar[rrr]^-{c^p_{Y'/X'}} & & &
(f')^*\Omega^p_{X'/\mathbf{Z}} \otimes_{\mathcal{O}_{Y'}} \det(\NL_{Y'/X'})
}
$$
commute for all $p \geq 0$. Here the vertical arrows use the maps
$b^*\Omega^p_{Y/\mathbf{Z}} \to \Omega^p_{Y'/\mathbf{Z}}$ and
$a^*\Omega^p_{X/\mathbf{Z}} \to \Omega^p_{X'/\mathbf{Z}}$
of Section \ref{section-de-rham-complex}
as well as the identification $b^*\det(\NL_{Y/X}) = \det(\NL_{Y'/X'})$ of
Discriminants, Section \ref{discriminant-section-tate-map}.
\end{lemma}

\begin{proof}
The map
$$
(f')^*\Omega_{X'/\mathbf{Z}} \oplus b^*\Omega_{Y/\mathbf{Z}}
\longrightarrow
\Omega_{Y'/\mathbf{Z}}
$$
is surjective because $Y'$ is an open subscheme of the fibre product;
the corresponding algebra statement is that $\Omega_{B \otimes_A C/\mathbf{Z}}$
is generated as a module by the images of $\Omega_{B/\mathbf{Z}}$ and
$\Omega_{C/\mathbf{Z}}$. Thus for all $p$ the map
$$
\bigoplus\nolimits_{i = 0, \ldots, p}
(f')^*\Omega^i_{X'/\mathbf{Z}} \otimes
b^*\Omega^{p - i}_{Y/\mathbf{Z}}
\longrightarrow
\Omega^p_{Y'/\mathbf{Z}}
$$
is surjective. Conditions (\ref{item-degree-zero}) and
(\ref{item-multiplicative}) combined with the commutativity
of the diagram in the lemma, implies that the map $c^p_{Y'/X'}$
is determined (but existence is not immediate).
\end{proof}

\begin{lemma}
\label{lemma-existence-base-change-Garel-upstairs}
Consider a commutative diagram
$$
\xymatrix{
Y' \ar[d]_{f'} \ar[r]_b & Y \ar[d]^f \\
X' \ar[r]^a & X
}
$$
of schemes which induces an isomorphism of $Y'$ with an open subscheme of
$X' \times_X Y$. Assume $f$ is locally quasi-finite and syntomic. Then for
every $y' \in Y'$ we can find opens $V' \subset Y'$,
$V \subset Y$, $U' \subset X$, $U \subset X$ with $y' \in V'$, with
$f'(V') \subset U'$, $b(V') \subset V$, $a(U') \subset U$,
$f(V) \subset U$, and such that for $p \geq 0$ there are maps
$c^p_{V/U}$ and $c^p_{V'/U'}$ satisfying (\ref{item-degree-zero})
and (\ref{item-multiplicative}) which are compatible with the diagram
$$
\xymatrix{
V' \ar[d] \ar[r]_b & V \ar[d] \\
U' \ar[r] & U
}
$$
in the sense explained in Lemma \ref{lemma-base-change-Garel-upstairs}.
\end{lemma}

\begin{proof}
It is clear that we may replace $Y'$ by $X' \times_X Y$ in order to prove this.
Pick affine opens $V \subset Y$ and $U \subset X$ with $V$ containing
the image of $y'$ as in
Discriminants, Lemma \ref{discriminant-lemma-syntomic-quasi-finite} part (5).
Pick an affine open $U' \subset X'$ containing the image of $y'$
and mapping into $U$. After replacing $X, Y, X', Y'$ by
$U, V, U', U' \times_U V$ we may assume we are in the situation
described in the next paragraph.

\medskip\noindent
Assume that $X' = \Spec(A')$, $X = \Spec(A)$, $Y = \Spec(B)$, and
$Y' = \Spec(A' \otimes_A B)$ with $B = A[x_1, \ldots, x_n]/(f_1, \ldots, f_n)$
a relative global complete intersection over $A$. Then also
$B' = A'[x_1, \ldots, x_n]/(f'_1, \ldots, f'_n)$ is a relative
global complete intersection over $A'$ where $f'_i$ is the image of $f_i$
in $A'[x_1, \ldots, x_n]$, see Algebra, Lemma
\ref{algebra-lemma-base-change-relative-global-complete-intersection}.
The construction in Remark \ref{remark-local-description}
provides us with the maps $c^p_{V/U}$ and $c^p_{V'/U'}$.
These maps are compatible by Lemma \ref{lemma-funny-map-functorial}
and the compatibility of the presentations of $B$ and $B'$ over
$A$ and $A'$. Some details omitted.
\end{proof}

\begin{proof}
This proof is very similar to the proof of
Discriminants, Proposition \ref{discriminant-proposition-tate-map}
and we suggest the reader look at that proof first.

\medskip\noindent
Let us reformulate the statement. Consider the category
$\mathcal{C}$ whose objects, denoted $Y/X$, are locally quasi-finite syntomic
morphism $f : Y \to X$ of schemes and whose morphisms
$b/a : Y'/X' \to Y/X$ are commutative diagrams
$$
\xymatrix{
Y' \ar[d]_{f'} \ar[r]_b & Y \ar[d]^f \\
X' \ar[r]^a & X
}
$$
which induce an isomorphism of $Y'$ with an open subscheme of
$X' \times_X Y$. The lemma means that for every object
$Y/X$ of $\mathcal{C}$ we have maps $c^p_{Y/X}$, $p \geq 0$
with properties (\ref{item-degree-zero}) and (\ref{item-multiplicative}),
and for every morphism $b/a : Y'/X' \to Y/X$ of $\mathcal{C}$ we have
that $c^p_{Y/X}$ and $c^p_{Y'/X'}$ are compatible as in
Lemma \ref{lemma-base-change-Garel-upstairs}.

\medskip\noindent
Given $Y/X$ in $\mathcal{C}$ and $y \in Y$ we can find
an affine open $V \subset Y$ and $U \subset X$ with $f(V) \subset U$
such that there exists some maps
$$
\Omega^p_{Y/\mathbf{Z}}|_V
\longrightarrow
\left(
f^*\Omega^p_{X/\mathbf{Z}} \otimes_{\mathcal{O}_Y} \det(\NL_{Y/X})
\right)|_V
$$
with properties (\ref{item-degree-zero}) and (\ref{item-multiplicative}).
This follows from picking affine opens as in
Discriminants, Lemma \ref{discriminant-lemma-syntomic-quasi-finite} part (5)
and then using the construction in Remark \ref{remark-local-description}.

\medskip\noindent
Note that the \'etale locus of $f$ is exactly the set of points where
$\delta(\NL_{Y/X})$ does not vanish, see discussion in
Discriminants, Section \ref{discriminant-section-tate-map}.
In particular
$f^*\Omega^p_{X/\mathbf{Z}} \to \Omega^p_{Y/\mathbf{Z}}$
becomes an isomorphism after inverting $\delta(\NL_{Y/X})$.
We conclude that if $\Omega^p_{X/\mathbf{Z}}$ is finite locally free and
the annihilator of $\delta(\NL_{Y/X})$ in $\mathcal{O}_Y$ is zero, then
the local maps constructed in the previous paragraph
are unique and automatically glue!

\medskip\noindent
Let $\mathcal{C}_{nice} \subset \mathcal{C}$ denote the full subcategory
of $Y/X$ such that
\begin{enumerate}
\item $\Omega_{X/\mathbf{Z}}$ is locally free, and
\item the annihilator of $\delta(\NL_{Y/X})$ in $\mathcal{O}_Y$ is zero.
\end{enumerate}
By the remarks in the previous paragraph, we see that for any
object $Y/X$ of $\mathcal{C}_{nice}$ we have unique maps
$c^p_{Y/X}$, $p \geq 0$ satisfying conditions (\ref{item-degree-zero})
and (\ref{item-multiplicative}). If $b/a : Y'/X' \to Y/X$
is a morphism of $\mathcal{C}_{nice}$, then the maps
$c^p_{Y/X}$ and $c^p_{Y'/X'}$ are compatible as in
Lemma \ref{lemma-base-change-Garel-upstairs}: namely,
locally there do exist compatible maps by
Lemma \ref{lemma-existence-base-change-Garel-upstairs}
and the uniqueness just mentioned shows these maps agree
with $c^p_{Y/X}$ and $c^p_{Y'/X'}$.
In other words, we have solved the problem
on the full subcategory $\mathcal{C}_{nice}$. For $Y/X$ in $\mathcal{C}_{nice}$
we continue to denote $c^p_{Y/X}$ the solution we've just found.

\medskip\noindent
In fact, more generally, suppose we have a morphism $b/a : Y'/X' \to Y/X$
of $\mathcal{C}$ such that $Y/X$ is an object of $\mathcal{C}_{nice}$.
The same argument, but this time using the uniqueness in
Lemma \ref{lemma-base-change-Garel-upstairs}, tells us there are
maps $c^p_{Y'/X'}$, $p \geq 0$  satisfying conditions (\ref{item-degree-zero})
and (\ref{item-multiplicative}) compatible with the already constructed
maps $c^p_{Y/X}$. Let us call this the base change
$(b/a)^*c^p_{Y/X}$.

\medskip\noindent
Consider morphisms
$$
Y_1/X_1 \xleftarrow{b_1/a_1} Y/X \xrightarrow{b_2/a_2} Y_2/X_2
$$
in $\mathcal{C}$ such that $Y_1/X_1$ and $Y_2/X_2$ are objects
of $\mathcal{C}_{nice}$.
{\bf Claim.} The two base changes $(b_i/a_i)^*c^p_{Y_i/X_i}$, $i = 1, 2$
are equal. We will first show that the claim implies the lemma
and then we will prove the claim.

\medskip\noindent
Let $d, n \geq 1$ and consider the locally
quasi-finite syntomic morphism $Y_{n, d} \to X_{n, d}$
constructed in Discriminants, Example
\ref{discriminant-example-universal-quasi-finite-syntomic}.
Then $Y_{n, d}$ and $Y_{n, d}$ are irreducible schemes of finite type and
smooth over $\mathbf{Z}$. Namely, $X_{n, d}$ is a spectrum of a
polynomial ring over $\mathbf{Z}$ and $Y_{n, d}$ is an open subscheme
of such. The morphism $Y_{n, d} \to X_{n, d}$ is locally quasi-finite syntomic
and \'etale over a dense open, see Discriminants, Lemma
\ref{discriminant-lemma-universal-quasi-finite-syntomic-etale}.
Thus $\delta(\NL_{Y_{n, d}/X_{n, d}})$ is nonzero: for example we have
the local description of $\delta(\NL_{Y/X})$ in
Discriminants, Remark \ref{discriminant-remark-local-description-delta}
and we have the local description of \'etale morphisms in
Morphisms, Lemma \ref{morphisms-lemma-etale-at-point} part (8).
Now a nonzero section of an invertible module over an irreducible
regular scheme has vanishing annihilator. Thus
$Y_{n, d}/X_{n, d}$ is an object of $\mathcal{C}_{nice}$.

\medskip\noindent
Let $Y/X$ be an arbitrary object of $\mathcal{C}$. Let $y \in Y$.
By Discriminants, Lemma \ref{discriminant-lemma-locally-comes-from-universal}
we can find $n, d \geq 1$ and morphisms
$$
Y/X \leftarrow V/U \xrightarrow{b/a} Y_{n, d}/X_{n, d}
$$
of $\mathcal{C}$ such that $V \subset Y$ and $U \subset X$ are open.
Then we have the base change $c^p_{V/U} = (b/a)^*c^p_{Y_{n, d}/X_{n, d}}$.
The claim guarantees these locally constructed maps $c^p_{V/U}$ glue!
Thus we get a well defined global maps $c^p_{Y/X}$ with properties
(\ref{item-degree-zero}) and (\ref{item-multiplicative}).
Finally, let $b/a : Y'/X' \to Y/X$ be an arbitray morphism of $\mathcal{C}$.
We have the maps $c^p_{Y'/X'}$, $p \geq 0$ and $c^p_{Y/X}$, $p \geq 0$
constructed in this paragraph. To check they are compatible as in
Lemma \ref{lemma-base-change-Garel-upstairs} we may work locally
on $Y'/X'$ and $Y/X$. Thus we may assume that there exists a morphism
$Y/X \to Y_{n, d}/X_{n, d}$. By construction, both the family
$c^p_{Y'/X'}$ and the family $c^p_{X/Y}$ are compatible with
the map towards $Y_{n, d}/X_{n, d}$. It follows easily from this
(and the uniqueness in Lemma \ref{lemma-base-change-Garel-upstairs})
that $c^p_{Y'/X'}$ is compatible with $c^p_{Y/X}$.
Thus it remains to prove the claim.

\medskip\noindent
In the rest of the proof we prove the claim. We may pick a point
$y \in Y$ and prove the maps agree in an open neighbourhood of $y$.
Thus we may replace $Y_1$, $Y_2$ by open neighbourhoods of the
image of $y$ in $Y_1$ and $Y_2$. Therefore we may assume
$Y, X, Y_1, X_1, Y_2, X_2$ are affine, say they are the spectra
of rings $B, A, B_1, A_1, B_2, A_2$. Picture
$$
\xymatrix{
B_1 \ar[r] & B & B_2 \ar[l] \\
A_1 \ar[u] \ar[r] & A \ar[u] & A_2 \ar[l] \ar[u]
}
$$
By assumption the spectrum of $B$ is an affine open of
both the spectrum of $A \otimes_{A_1} B_1$ and $A \otimes_{A_2} B_2$.
Shrinking more we may assume there exist elements
$g_i \in A \otimes_{A_i} B_i$ such that our maps give isomorphisms
$(A \otimes_{A_i} B_i)_{g_i} = B$, see
Properties, Lemma \ref{properties-lemma-standard-open-two-affines}.
Let $x_\alpha$, $y_\beta$ be a sufficiently large collection of
variables such that we may choose surjections
$$
A'_1 = A_1[x_\alpha] \to A
\quad\text{and}\quad
A'_2 = A_2[y_\beta] \to A
$$
of $A_1$ and $A_2$-algebras. Then we can choose lifts
$h_i \in A'_i \otimes_{A_i} B_i$ of $g_i \in A \otimes_{A_i} B_i$
and we consider the diagram
$$
\xymatrix{
(A'_1 \otimes_{A_1} B_1)_{h_1} \ar[r] & B &
(A'_2 \otimes_{A_1} B_2)_{h_2} \ar[l] \\
A'_1 \ar[u] \ar[r] & A \ar[u] & A'_2 \ar[l] \ar[u]
}
$$
By construction the two squares are cocartesian. Next, we consider the
ring map
$$
A' = A'_1 \times_A A'_2 \longrightarrow
B' =
(A'_1 \otimes_{A_1} B_1)_{h_1} \times_B
(A'_2 \otimes_{A_1} B_2)_{h_2}
$$
By
More on Algebra, Lemma \ref{more-algebra-lemma-module-over-fibre-product}
we have $A'_1 \otimes_{A'} B' = (A'_1 \otimes_{A_1} B_1)_{h_1}$
and $A'_2 \otimes_{A'} B' = (A'_2 \otimes_{A_2} B_2)_{h_2}$. In particular
the fibres of the morphism $Y' = \Spec(B') \to \Spec(A') = X'$ are
open subschemes of base changes of the fibres of the maps $Y_i \to X_i$, hence
finite and local complete intersections (see
Discriminants, Lemma \ref{discriminant-lemma-syntomic-quasi-finite}).
By More on Algebra, Lemma
\ref{more-algebra-lemma-properties-algebras-over-fibre-product}
the ring map $A' \to B'$ is flat and of finite presentation.
Thus by Discriminants, Lemma
\ref{discriminant-lemma-syntomic-quasi-finite} part (3)
we conclude that $Y'/X'$ is an object of $\mathcal{C}$.
Consider now the commutative diagram
$$
\xymatrix{
& & Y/X \ar[ld] \ar@/_2pc/[lldd]_{b_1/a_1} \ar[rd] \ar@/^2pc/[rrdd]^{b_2/a_2} \\
& Y'_1/X'_1 \ar[rd] \ar[ld] & & Y'_2/X'_2 \ar[ld] \ar[rd] \\
Y_1/X_1 & & Y'/X' & & Y_2/X_2
}
$$
where $Y'_i/X'_i$ corresponds to $A'_i \to (A'_i \otimes_{A_i} B_i)_{h_i}$.
Note that $Y'_i/X'_i$ is an object of $\mathcal{C}_{nice}$ as it
obtained by taking an open subscheme of the product of an
infinite dimensional affine space with $Y_i/X_i$; small detail omitted.
In particular, the pullback of $c^p_{Y_i/X_i}$ via $(b_i/a_i)$ is the same
as the pullback of $c^p_{Y'_i/X'_i}$ via $Y/X \to Y'_i/X'_i$.
Now we would be done if $Y'/X'$ is an object of $\mathcal{C}_{nice}$,
but this is almost never the case. Namely, then pulling back $c^p_{Y'/X'}$
around the two sides of the square, we would obtain the desired conclusion.
To get around the problem that $Y'/X'$ is not in $\mathcal{C}_{nice}$
we note the arguments above show that, after possibly shrinking all
of the schemes $X, Y, X'_1, Y'_1, X'_2, Y'_2, X', Y'$ we can find some
$n, d \geq 1$, and extend the diagram like so:
$$
\xymatrix{
& Y/X \ar[ld] \ar[rd] \\
Y'_1/X'_1 \ar[rd] & & Y'_2/X'_2 \ar[ld] \\
& Y'/X' \ar[d] \\
& Y_{n, d}/X_{n, d}
}
$$
and then we can use the already given argument by pulling
back from $c^p_{Y_{n, d}/X_{n, d}}$. This finishes the proof.
\end{proof}

\begin{lemma}
\label{lemma-syntomic-quasi-finite}
Let $f : Y \to X$ be a morphism of schemes. The following are equivalent
\begin{enumerate}
\item $f$ is locally quasi-finite and syntomic,
\item $f$ is locally quasi-finite, flat, and a local complete intersection
morphism,
\item $f$ is locally quasi-finite, flat, locally of finite presentation,
and the fibres of $f$ are local complete intersections,
\item $f$ is locally quasi-finite and for every $y \in Y$ there are
affine opens $y \in V = \Spec(B) \subset Y$, $U = \Spec(A) \subset X$
with $f(V) \subset U$ an integer $n$ and
$h, f_1, \ldots, f_n \in A[x_1, \ldots, x_n]$ such that
$B = A[x_1, \ldots, x_n, 1/h]/(f_1, \ldots, f_n)$,
\item for every $y \in Y$ there are affine opens
$y \in V = \Spec(B) \subset Y$, $U = \Spec(A) \subset X$
with $f(V) \subset U$ such that $A \to B$ is a relative global complete
intersection of the form $B = A[x_1, \ldots, x_n]/(f_1, \ldots, f_n)$,
\item $f$ is locally quasi-finite, flat, locally of finite presentation,
and $\NL_{Y/X}$ has tor-amplitude in $[-1, 0]$, and
\item $f$ is flat, locally of finite presentation,
$\NL_{Y/X}$ is perfect of rank $0$ with tor-amplitude in $[-1, 0]$,
\end{enumerate}
\end{lemma}

\begin{proof}
The equivalence of (1) and (2) is
More on Morphisms, Lemma \ref{more-morphisms-lemma-flat-lci}.
The equivalence of (1) and (3) is 
Morphisms, Lemma \ref{morphisms-lemma-syntomic-flat-fibres}.

\medskip\noindent
If $A \to B$ is as in (4), then
$B = A[x, x_1, \ldots, x_n]/(xh - 1, f_1, \ldots, f_n)$
is a relative global complete intersection by see Algebra, Definition
\ref{algebra-definition-relative-global-complete-intersection}.
Thus (4) implies (5).
It is clear that (5) implies (4).

\medskip\noindent
Condition (5) implies (1): by
Algebra, Lemma \ref{algebra-lemma-relative-global-complete-intersection}
a relative global complete intersection is syntomic and
the definition of a relative global complete intersection
guarantees that a relative global complete intersection on
$n$ variables with $n$ equations is quasi-finite, see
Algebra, Definition
\ref{algebra-definition-relative-global-complete-intersection} and
Lemma \ref{algebra-lemma-isolated-point-fibre}.

\medskip\noindent
Either Algebra, Lemma \ref{algebra-lemma-syntomic} or
Morphisms, Lemma \ref{morphisms-lemma-syntomic-locally-standard-syntomic}
shows that (1) implies (5).

\medskip\noindent
More on Morphisms, Lemma \ref{more-morphisms-lemma-flat-fp-NL-lci} shows that
(6) is equivalent to (1). If the equivalent conditions (1) -- (6) hold,
then we see that affine locally $Y \to X$ is given by a relative global
complete intersection $B = A[x_1, \ldots, x_n]/(f_1, \ldots, f_n)$
with the same number of variables as the number of
equations. Using this presentation we see that
$$
\NL_{B/A} =\left(
(f_1, \ldots, f_n)/(f_1, \ldots, f_n)^2
\longrightarrow
\bigoplus\nolimits_{i = 1, \ldots, n} B \text{d} x_i\right)
$$
By Algebra, Lemma
\ref{algebra-lemma-relative-global-complete-intersection-conormal}
the module $(f_1, \ldots, f_n)/(f_1, \ldots, f_n)^2$
is free with generators the congruence classes of the elements
$f_1, \ldots, f_n$. Thus $\NL_{B/A}$ has rank $0$ and so does $\NL_{Y/X}$.
In this way we see that (1) -- (6) imply (7).

\medskip\noindent
Finally, assume (7). By
More on Morphisms, Lemma \ref{more-morphisms-lemma-flat-fp-NL-lci}
we see that $f$ is syntomic. Thus on suitable affine opens
$f$ is given by a relative global complete intersection
$A \to B = A[x_1, \ldots, x_n]/(f_1, \ldots, f_m)$, see
Morphisms, Lemma \ref{morphisms-lemma-syntomic-locally-standard-syntomic}.
Exactly as above we see that $\NL_{B/A}$ is a perfect complex
of rank $n - m$. Thus $n = m$ and we see that (5) holds.
This finishes the proof.
\end{proof}

\begin{lemma}
\label{lemma-characterize-invertible}
Invertibility of the relative dualizing module.
\begin{enumerate}
\item If $A \to B$ is a quasi-finite flat homomorphism of Noetherian rings,
then $\omega_{B/A}$ is an invertible $B$-module if and only if
$\omega_{B \otimes_A \kappa(\mathfrak p)/\kappa(\mathfrak p)}$
is an invertible $B \otimes_A \kappa(\mathfrak p)$-module
for all primes $\mathfrak p \subset A$.
\item If $Y \to X$ is a quasi-finite flat morphism of
Noetherian schemes, then $\omega_{Y/X}$ is invertible
if and only if $\omega_{Y_x/x}$ is invertible for all $x \in X$.
\end{enumerate}
\end{lemma}

\begin{proof}
Proof of (1). As $A \to B$ is flat, the module
$\omega_{B/A}$ is $A$-flat, see Lemma \ref{lemma-dualizing-base-flat-flat}.
Thus $\omega_{B/A}$ is an invertible $B$-module if and only if
$\omega_{B/A} \otimes_A \kappa(\mathfrak p)$
is an invertible $B \otimes_A \kappa(\mathfrak p)$-module for
every prime $\mathfrak p \subset A$, see More on Morphisms, Lemma
\ref{more-morphisms-lemma-flat-and-free-at-point-fibre}.
Still using that $A \to B$ is flat, we have that
formation of $\omega_{B/A}$ commutes with base change, see
Lemma \ref{lemma-dualizing-base-change-of-flat}.
Thus we see that invertibility of the relative dualizing module,
in the presence of flatness, is equivalent to invertibility
of the relative dualizing module for the maps
$\kappa(\mathfrak p) \to B \otimes_A \kappa(\mathfrak p)$.

\medskip\noindent
Part (2) follows from (1) and the fact that affine locally
the dualizing modules are given by their algebraic counterparts, see
Remark \ref{remark-relative-dualizing-for-quasi-finite}.
\end{proof}

\begin{lemma}
\label{lemma-dim-zero-global-complete-intersection-over-field}
Let $k$ be a field. Let $B = k[x_1, \ldots, x_n]/(f_1, \ldots, f_n)$
be a global complete intersection over $k$ of dimension $0$.
Then $\omega_{B/k}$ is invertible.
\end{lemma}

\begin{proof}
By Noether normalization, see
Algebra, Lemma \ref{algebra-lemma-Noether-normalization}
we see that there exists a finite injection $k \to B$, i.e.,
$\dim_k(B) < \infty$. Hence $\omega_{B/k} = \Hom_k(B, k)$
as a $B$-module.
By Dualizing Complexes, Lemma \ref{dualizing-lemma-dualizing-finite}
we see that $R\Hom(B, k)$ is a dualizing complex for $B$
and by Dualizing Complexes, Lemma \ref{dualizing-lemma-RHom-ext}
we see that $R\Hom(B, k)$ is equal to $\omega_{B/k}$
placed in degree $0$. Thus it suffices to show that
$B$ is Gorenstein
(Dualizing Complexes, Lemma \ref{dualizing-lemma-gorenstein}).
This is true by Dualizing Complexes, Lemma
\ref{dualizing-lemma-gorenstein-lci}.
\end{proof}

\begin{lemma}
\label{lemma-dualizing-syntomic-quasi-finite}
Let $f : Y \to X$ be a morphism of locally Noetherian schemes. If $f$
satisfies the equivalent conditions of Lemma \ref{lemma-syntomic-quasi-finite}
then $\omega_{Y/X}$ is an invertible $\mathcal{O}_Y$-module.
\end{lemma}

\begin{proof}
We may assume $A \to B$ is a relative global complete
intersection of the form $B = A[x_1, \ldots, x_n]/(f_1, \ldots, f_n)$
and we have to show $\omega_{B/A}$ is invertible.
This follows in combining Lemmas \ref{lemma-characterize-invertible} and
\ref{lemma-dim-zero-global-complete-intersection-over-field}.
\end{proof}

\begin{example}
\label{example-universal-quasi-finite-syntomic}
Let $n \geq 1$ and $d \geq 1$ be integers. Let $T$ be the set of
multi-indices $E = (e_1, \ldots, e_n)$ with $e_i \geq 0$ and
$\sum e_i \leq d$. Consider the ring
$$
A = \mathbf{Z}[a_{i, E} ; 1 \leq i \leq n, E \in T]
$$
In $A[x_1, \ldots, x_n]$ consider the elements
$f_i = \sum_{E \in T} a_{i, E} x^E$ where $x^E = x_1^{e_1} \ldots x_n^{e_n}$
as is customary. Consider the $A$-algebra
$$
B = A[x_1, \ldots, x_n]/(f_1, \ldots, f_n)
$$
Denote $X_{n, d} = \Spec(A)$ and let $Y_{n, d} \subset \Spec(B)$
be the maximal open subscheme such that the restriction of the
morphism $\Spec(B) \to \Spec(A) = X_{n, d}$ is quasi-finite, see
Algebra, Lemma \ref{algebra-lemma-quasi-finite-open}.
\end{example}

\begin{lemma}
\label{lemma-universal-quasi-finite-syntomic-etale}
With notation as in Example \ref{example-universal-quasi-finite-syntomic}
the schemes $X_{n, d}$ and $Y_{n, d}$ are regular and irreducible,
the morphism $Y_{n, d} \to X_{n, d}$ is locally quasi-finite and
syntomic, and there is a dense open subscheme $V \subset Y_{n, d}$
such that $Y_{n, d} \to X_{n, d}$ restricts to an \'etale morphism
$V \to X_{n, d}$.
\end{lemma}

\begin{proof}
The scheme $X_{n, d}$ is the spectrum of the polynomial ring $A$.
Hence $X_{n, d}$ is regular and irreducible. Since we can write
$$
f_i = a_{i, (0, \ldots, 0)} +
\sum\nolimits_{E \in T, E \not = (0, \ldots, 0)} a_{i, E} x^E
$$
we see that the ring $B$ is isomorphic to the polynomial ring
on $x_1, \ldots, x_n$ and the elements $a_{i, E}$ with
$E \not = (0, \ldots, 0)$. Hence $\Spec(B)$ is an irreducible and
regular scheme and so is the open $Y_{n, d}$. The morphism
$Y_{n, d} \to X_{n, d}$ is locally quasi-finite and syntomic by
Lemma \ref{lemma-syntomic-quasi-finite}. To find $V$ it suffices
to find a single point where $Y_{n, d} \to X_{n, d}$ is \'etale
(the locus of points where a morphism is \'etale is open by
definition). Thus it suffices to find a point of $X_{n, d}$
where the fibre of $Y_{n, d} \to X_{n, d}$ is nonempty and \'etale, see
Morphisms, Lemma \ref{morphisms-lemma-etale-at-point}. We choose
the point corresponding to the ring map $\chi : A \to \mathbf{Q}$
sending $f_i$ to $x_i^d - 1$. Then
$$
B \otimes_{A, \chi} \mathbf{Q} =
\mathbf{Q}[x_1, \ldots, x_n]/(x_1^d - 1, \ldots, x_n^d - 1)
$$
which is a nonzero \'etale algebra over $\mathbf{Q}$.
\end{proof}

\begin{lemma}
\label{lemma-locally-comes-from-universal}
Let $f : Y \to X$ be a morphism of schemes. If $f$ satisfies the equivalent
conditions of Lemma \ref{lemma-syntomic-quasi-finite} then for every
$y \in Y$ there exist $n, d$ and a commutative diagram
$$
\xymatrix{
Y \ar[d] &
V \ar[d] \ar[l] \ar[r] &
Y_{n, d} \ar[d] \\
X & U \ar[l] \ar[r] &
X_{n, d}
}
$$
where $U \subset X$ and $V \subset Y$ are open with $y \in V$,
where $Y_{n, d} \to X_{n, d}$
is as in Example \ref{example-universal-quasi-finite-syntomic}, and
where the square on the right hand side is cartesian.
\end{lemma}

\begin{proof}
By Lemma \ref{lemma-syntomic-quasi-finite}
we can choose $U$ and $V$ affine so that
$U = \Spec(R)$ and $V = \Spec(S)$ with
$S = R[y_1, \ldots, y_n]/(g_1, \ldots, g_n)$.
With notation as in Example \ref{example-universal-quasi-finite-syntomic}
if we pick $d$ large enough, then we can write each $g_i$ as
$g_i = \sum_{E \in T} g_{i, E}y^E$ with $g_{i, E} \in R$.
Then the map $A \to R$ sending $a_{i, E}$ to $g_{i, E}$
and the map $B \to S$ sending $x_i \to y_i$ give a cocartesian
diagram of rings
$$
\xymatrix{
S & B \ar[l] \\
R \ar[u] & A \ar[l] \ar[u]
}
$$
which proves the lemma.
\end{proof}

\bibliographystyle{plain}
\bibliography{my}
\end{document}
